%قبل از شروع مسئله ترجیحاً چیزی ننویسید. اما اگر خواستید بنویسید، در قالب لیستی (فقرات و شمارش) بنویسید.
%برای اضافه کردن نمره، بعدی از شروع‌مسئله، [x] را اضافه کنید و به جای x  نمره مورد نظر را قرار دهید. برای زیرمسئله‌ها هم به همین صورت است.
%\شروع‌مسئله[10]
\شروع‌مسئله
%\متن‌سیاه{سوال عملی اول:}\\
%\متن‌سیاه{دراپ‌تست}\\
در این چالش شما باید یک برنامه‌ی فشرده‌سازی برای متون فارسی بنویسید. هدف از این چالش آشنایی شما با الگوریتم‌های فشرده‌سازی و استفاده از داده‌ساختارهای مناسب برای بهبود عملیات فشرده‌سازی است. برای آشنایی با الگوریتم‌های مقدماتی فشرده‌سازی می‌توانید از 
\href{http://www.ijarcst.com/doc/vol2-issue3/ver.3/hardik_jani.pdf}{این لینک}
استفاده کنید. دقت کنید که بسیاری از این الگوریتم‌ها در سطح کاراکتر فشرده‌سازی می‌کنند در حالی که برای رسیدن به مناسب‌ترین فشرده‌سازی برای متون باید از کلمات در فشرده‌سازی استفاده کرد. شما باید یک برنامه برای فشرده‌سازی (\متن‌لاتین{compressor}) متن و یک برنامه برای برگرداندن آن به حالت اولیه (\متن‌لاتین{decompressor}) بنویسید و آن‌ها را به فرمتی که گفته می‌شود در کوئرا آپلود کنید. پیشنهاد می‌شود جست‌و‌جوی الگوریتم‌های مختلف فشرده‌سازی را به روزهای آخر موکول نکنید چرا که فرصت پیاده‌سازی، تست و بهبود آن‌ها را پیدا نخواهید کرد.

\متن‌سیاه{ورودی و خروجی:}

ورودی برنامۀ فشرده‌سازی شما یک فایل شامل متن فارسی با نام \متن‌لاتین{input.txt} است و خروجی آن فایل فشرده‌سازی شده با نام \متن‌لاتین{output.txt} است. به همین ترتیب ورودی برنامۀ بازگردانی شما یک فایل فشرده‌سازی شده با نام \متن‌لاتین{output.txt} است و خروجی آن همان فایل اولیه متن فارسی با نام \متن‌لاتین{input.txt} است. فایل‌های ورودی و خروجی در شاخۀ جاری اجرای برنامه قرار دارند. دقت کنید که ورودی فشرده‌سازی و خروجی بازگردانی باید کاملاً یکسان باشند، در غیر این صورت نمره‌ای به شما تعلق نمی‌گیرد. ممکن است لازم باشد برنامۀ شما فایل \متن‌لاتین{input.txt} را به صورت باینری (و نه تکست) بخواند. همچنین دقت کنید که حجم کد شما حداکثر $1$ مگابایت می‌تواند باشد.

\متن‌سیاه{نحوۀ نمره‌دهی:}

برنامۀ شما از دو جنبۀ زمان اجرا و میزان فشرده‌سازی بررسی خواهد شد. زمان اجرای برنامۀ شما برابر مجموع زمان فشرده‌سازی و زمان  بازگردانی است. برنامۀ شما بر روی چند فایل متن فارسی اجرا خواهد شد و امتیازتان برای هر فایل محاسبه خواهد شد. امتیاز شما برای یک فایل ترکیبی خطی از میزان فشرده‌سازی و زمان اجرا است که تاثیر میزان فشرده‌سازی در آن بیشتر است. نحوۀ دقیق این امتیازدهی متعاقباً اعلام خواهد شد. در نهایت امتیاز نهایی شما مجموع امتیازات کد شما بر روی فایل‌های مختلف است. افراد بر اساس امتیازهایشان رده‌بندی می‌شوند و نمره‌دهی بر اساس این رده‌بندی به صورت زیر انجام می‌شود: به نفر اول نمرۀ $1.25$ تعلق می‌گیرد و نمرۀ بقیۀ افراد به صورت خطی کاهش پیدا می‌کند. همچنین در صورتی که برنامۀ شما به درستی فشرده‌سازی و بازگردانی را (در زمان معقولی) انجام دهد و نرخ فشرده‌سازی شما کمتر از $75$ درصد باشد، حداقل $0.5$ نمره به شما تعلق خواهد گرفت.
%\متن‌سیاه{ورودی:}\\
%در خط اول ورودی، ابتدا k و سپس m داده می‌شود.\\
%\متن‌سیاه{خروجی:}\\
%در تنها سطر خروجی، تعداد آیفون‌های مورد نیاز می‌آید. اگر در هیچ صورتی چنین کاری ممکن نبود ۱- چاپ می‌شود.\\
%\متن‌سیاه{محدودیت‌ها:}\\
%$1\leq k\leq 10^4$

%$m\leq10^9$
%\\
%\textbf{\متن‌سیاه{ورودی و خروجی نمونه}\\
%% تست اول
%\begin{latin}
%	\begin{tabularx}{\textwidth}{|X | X|}
%		\hline
%		\textbf{stdin} &
%		\textbf{stdout}\\
%		\hline
%		100 100&
%		1\\
%		\hline
%	
%	\end{tabularx}
%\end{latin}
%
%% تست دوم
%
%\begin{latin}
%	\begin{tabularx}{\textwidth}{|X | X|}
%		\hline
%		\textbf{stdin} &
%		\textbf{stdout}\\
%		\hline
%		340 20&
%		3\\
%		\hline
%		
%	\end{tabularx}
%\end{latin}}
% صورت اصلی مسئله را در این‌جا بنویسید.
%اگر مسئله چند قسمت داشت به این صورت عمل کنید:
%\شروع{زیرمسائل}
%\زیرمسئله[0] % مسئلۀ ۱آ
%\زیرمسئله % مسئلۀ ۱ب
%\زیرمسئله % مسئلۀ ۱پ
%\زیرمسئله % مسئلۀ ۱ت
%\پایان{زیرمسائل}



%%%%%%%%%%%%%%%%%%%%%%%%%%%%%%%%%%%%%%%%%%%%%%%%%%%%%
% برخی دستورات مفید برای تمرین‌ها. %
%%%%%%%%%%%%%%%%%%%%%%%%%%%%%%%%%%%%%%%%%%%%%%%%%%%%%

% فرمول‌های نمونه:
%$f(x) = \Theta(x)$
%$T(x, y) \leq \log(x) + 2^y + \binom{2n}{n}$
% {\tt code\_function}


% لیست موارد بدون شماره‌گذاری:
%\شروع{فقرات}
%    \فقره مورد اول 
%        \شروع{فقرات}
%            \فقره مورد اول داخلی 
%                \شروع{فقرات}
%                    \فقره اولین مورد خیلی داخلی 
%                    \فقره دومین مورد خیلی داخلی 
%                \پایان{فقرات}
%            \فقره مورد دوم داخلی 
%        \پایان{فقرات}
%    \فقره مورد دوم 
%\پایان{فقرات}

% لیست موارد شماره‌گذاری شده:
%\شروع{شمارش}
%    \فقره یک مورد 
%    \فقره یک مورد دیگر 
%    \فقره باز هم یک مورد دیگر
%\پایان{شمارش}

% برای نوشتن یک فرمول چند خطی که یک شماره بخوره. برای حالت‌های دیگر به سایت زیر مراجعه کنید:
% https://www.overleaf.com/learn/latex/Aligning_equations_with_amsmath
%\begin{align} 
%    \begin{split}
%        (x+y)^3 &= (x+y)^2(x+y) \\
%                &= (x^2+2xy+y^2)(x+y) \\
%                &= (x^3+2x^2y+xy^2) + (x^2y+2xy^2+y^3) \\
%                &= x^3+3x^2y+3xy^2+y^3
%    \end{split}                                 
%\end{align}

% نوشتن ماتریس و آرایه دوبعدی
%\begin{align}
%    A = 
%    \begin{bmatrix}
%        A_{11} & A_{21} \\
%        A_{21} & A_{22}
%    \end{bmatrix}
%\end{align}

% اضافه کردن تصویر با شرح. برای اضافه کردن تصویر بدون شرح، فقط دستور درج‌تصویر را استفاده کنید.
%\شروع{شکل}[h!]
%\تنظیم‌ازوسط
%  \درج‌تصویر[پهنا=0.5\پهنای‌متن]{img.jpg}
%\شرح{شرح تصویر}
%\برچسب{تصویر:۱}
%\پایان{شکل}
%\پایان{زیرمسائل}

%کد زیر را ببینید.
%%\lstinputlisting{source_filename.py}
%\begin{latin}
%\begin{lstlisting}
%# Import the modules
%import sys
%import random
%
%ans = True
%
%while ans:
%question = raw_input("Ask the magic 8 ball a question: (press enter to quit) ")
%
%answers = random.randint(1,8)
%
%if question == "":
%sys.exit()
%
%elif answers == 1:
%print "It is certain"
%
%elif answers == 2:
%print "Outlook good"
%
%elif answers == 3:
%print "You may rely on it"
%
%elif answers == 4:
%print "Ask again later"
%
%elif answers == 5:
%print "Concentrate and ask again"
%
%elif answers == 6:
%print "Reply hazy, try again"
%
%elif answers == 7:
%print "My reply is no"
%
%elif answers == 8:
%print "My sources say no"
%\end{lstlisting}
%\end{latin}
